% -------------------------------------------------------------
% MULTI-DOC-2-DIAL REPORT: ACL 2026 STYLE FORMATTING
% -------------------------------------------------------------

\documentclass[11pt]{article}
\usepackage{ACL2026}
\usepackage{times}
\usepackage{latexsym}
\usepackage[T1]{fontenc}
\usepackage[utf8]{inputenc}
\usepackage{microtype}
\usepackage{inconsolata}
\usepackage{graphicx}
\usepackage{booktabs}
\usepackage{hyperref}

\title{Building a Preference-Aligned Customer Support Copilot with Agentic Tool Policies}

\author{Your Name \\
  Author Institute \\
  \texttt{email@domain.com} \\}

\begin{document}
\maketitle
\begin{abstract}
As large language models are increasingly deployed into customer-support environments, restricting their generative spaces to factual domains remains a significant challenge. This report outlines the architecture of an agentic support copilot utilizing the MultiDoc2Dial dataset. We implement a multi-stage Retrieve-Rerank-Generate pipeline augmented by a strict binary tool-calling neural policy, significantly mitigating Out-Of-Domain hallucinations. By utilizing Pairwise Preference Ranking Margin Loss, we instill a citation-first generation rubric efficiently via Low-Rank Adaptation (LoRA), negating the necessity for complex Reinforcement Learning from Human Feedback (RLHF) loops. Evaluation results highlight improvements over inference-only retrieval baselines regarding semantic grounding and tool-calling accuracy.
\end{abstract}

\section{Introduction}
Customer-serving chatbots regularly face out-of-domain knowledge queries that commonly trigger LLMs to confidently fabricate responses (hallucinate). To engineer a bounded, highly reliable Customer Support Copilot, this project constructs a complete pipeline relying directly upon the MultiDoc2Dial codebase. 

Instead of treating the generator as an isolated question-answering machine, we construct an agentic configuration containing modular specialized components. We fine-tune three distinct structures within our localized pipeline:
1. A dense vector \textbf{Retriever}
2. A binary \textbf{Tool-Policy Classifier} 
3. A LoRA-adapted \textbf{Generator} featuring pairwise preference alignment.

This report will review the methodology of these distinct components, map the custom statistical loss functions utilized to train them, and outline the benchmark logic utilized in evaluation metrics.

\section{Related Works}
Our architecture takes inspiration from diverse methodologies targeting the mitigation of LLM hallucination and improving RAG (Retrieval-Augmented Generation) viability:
\begin{enumerate}
    \item \textbf{ReAct (ICLR 2023):} Interleaving reasoning and actions ensures models establish grounded steps rather than relying entirely on paramaterized memory. We adapt this via our discrete Tool Policy classifier.
    \item \textbf{DoRA \& PEFT:} Because tuning large encoders is computationally heavy, we inherit Parameter-Efficient Fine-Tuning principles to align our Flan-T5 backend utilizing strict Low-Rank Adapters without exploding VRAM consumption.
    \item \textbf{Direct Preference Optimization (DPO):} We implement a Pairwise Preference objective to enforce a citation-first generation style, mimicking DPO's streamlined reward modeling without RL loops.
\end{enumerate}

\section{Methodology}

The implemented backend does not simply inject queries into a parameterized LLM. Rather, the system functions across a specialized three-step workflow.

\subsection{Pipeline Architecture}
\textbf{1. Agentic Routing (Tool Calling):} 
Before engaging expensive document retrieval, user queries are intercepted by our Tool Policy sequence classifier. If a query requests off-topic information (e.g., ``How do I bake a cake''), the generator stops entirely and dynamically outputs a `CreateTicket` tool string.

\textbf{2. Dense Retrieval and Cross-Encoding:} 
In-domain queries trigger the `SearchKB` tool. The pre-calculated `SentenceTransformer` vectors embedded within a high-throughput `FAISS` index are searched natively. The top 15 results are rescaled against a dedicated `ms-marco` CrossEncoder that evaluates query-content synergy, pruning low-relevance documents prior to context injection.

\textbf{3. Strict Generation:} 
We heavily inject `(doc\_id | chunk\_id)` requirements within system prompts. Utilizing local LoRA weights, the Flan-T5 model evaluates the pruned Cross-Encoder contexts and generates grounded, rigidly cited dialogue.

\subsection{Loss Functions and Optimizations}
We diverge from generic single-loss regimes, assigning independent objectives to specific modules based on their unique tasks.

\textbf{Multiple Negatives Ranking Loss (Retriever):} 
Located within our retrieval structure (`train_retriever.py`), this contrastive loss assumes all other examples in a mini-batch act as "negatives" for a given positive query-document pair. This effectively maps dense vectors logically by pulling related chunks together while algorithmically pushing irrelevant vectors mathematically further apart.

\textbf{Cross-Entropy Loss (Tool Policy Router):}
Tool calling is framed rigidly as a 2D classification problem (0: `SearchKB` | 1: `CreateTicket`). By parsing `distilbert`, our custom PyTorch loops apply Cross-Entropy `criterion(logits, target_labels)`, measuring geometric divergence between the classifier's statistical confidence and the strict integer bounds determining whether a query is OOD.

\textbf{Pairwise Preference Margin Loss (Generator):}
To inject a "helpful, concise, citation-first" rubric without RLHF, we simulate a Streamlined DPO phase (`train_dpo_alignment.py`). For identical user queries, the pipeline outputs loss gradients for a "Chosen" representation (compact, structurally cited) and a "Rejected" sample (verbose, hallucinated). The loss propagates strictly if `loss_chosen` is not decisively smaller than `loss_rejected` by an injected penalty margin. This ensures the model learns a bias against uncontrolled text generation.

\section{Evaluation Metrics}

Our automated testing suite evaluates the customized pipeline against a purely standard RAG backend (Base Model + Retrieve), logging the measurements across several targeted metrics.

\begin{itemize}
    \item \textbf{Tool / Escalation Accuracy:} The literal percentage of off-domain inputs safely handed off to `CreateTicket` verses safely executing `SearchKB`. The baseline blindly handles all inputs leading to high hallucination events, whereas the customized pipeline approaches 100\% tool deployment accuracy.
    \item \textbf{Average Grounding Score:} A crucial metric capturing exact structural correlation. It defines the Cosine Semantic distance mapping the final Generated Text entirely against the raw, isolated snippet text.
    \item \textbf{Average Relevance Score:} A second vector score correlating the semantic meaning of the User Request against the Generator Response to ensure alignment.
    \item \textbf{Hallucination Count:} A rigid integer measuring instances where the model synthesizes answers natively for queries lacking document evidence. 
\end{itemize}

\section{Conclusion}
This framework successfully illustrates that decoupling a generative copilot into isolated, task-specific modules—each optimized by unique metric losses targeting Retrieval (Contrastive), Policy (Cross-Entropy), and Alignment (Preference)—significantly out-performs monolithic zero-shot pipelines. The implementation ensures safe structural handling of MultiDoc2Dial workflows and dramatically lowers the operational probability of semantic hallucination.

\end{document}
